\chapter{微分几何与代数拓扑基础速查}
\label{sec:section_4647}

\textbf{(Appendix A: Cheat Sheet for Differential Geometry \& Algebraic Topology)}

本节为工程师提供必要的数学词汇表，将 MSC 的概念映射到标准的数学定义上。

\smalltitle{纤维丛 (Fiber Bundle) — $\mathcal{U} = (E, \pi, M, F)$}
\begin{itemize}
\item   \textbf{数学定义}：一个局部同胚于 $\mathcal{M} \times F$ 的空间 $E$。

\item   \textbf{MSC 对应}：\textbf{宇宙/知识库的本体论结构}。

\item   $\mathcal{M}$ (底空间)：\textbf{形 (Morphos)}。拓扑骨架，回答 "Where/How connected"。

\item   $F$ (纤维)：\textbf{质 (Qualia)}。属性向量空间，回答 "What"。

\item   $E$ (全空间)：\textbf{实存 (Existence)}。形与质结合后的状态空间。

\item   $\pi$ (投影)：\textbf{抽象 (Abstraction)}。忽略属性，只看结构的操作。
\end{itemize}



\smalltitle{截面 (Section) — $\sigma:\mathcal{M}\to E$}

\begin{itemize}
\item   \textbf{数学定义}：满足 $\pi \circ \sigma = \text{id}_M$ 的连续映射。

\item   \textbf{MSC 对应}：\textbf{实体 (Entity) / 现象 (Phenomenon)}。

\item   一个物体就是纤维丛上的一个\textbf{波形切片}。

\item   \textit{全局截面}：完整的世界状态。

\item   \textit{局部截面}：单一物体或局部场。
\end{itemize}



\smalltitle{联络 (Connection) — $\nabla$}

\begin{itemize}
\item   \textbf{数学定义}：定义在切丛上的微分算子，用于区分水平方向（形变）和垂直方向（质变）。

\item   \textbf{MSC 对应}：\textbf{语义的一致性 / 物理定律}。

\item   \textbf{平行移动 (Parallel Transport)}：当物体移动时，如何保持其“质”不变（物体恒常性）。

\item   \textbf{规范场 (Gauge Field)}：引起“质”发生相位旋转的力（如电磁力/情感偏置）。
\end{itemize}



\smalltitle{曲率 (Curvature) — $\Omega = d\omega + \omega \wedge \omega$}

\begin{itemize}
\item   \textbf{数学定义}：沿闭合路径平移一圈后的偏差。

\item   \textbf{MSC 对应}：\textbf{力 (Force) / 惊奇 (Surprisal)}。

\item   在物理中，曲率是引力。

\item   在语义中，曲率是\textbf{语境张力}。如果一个词在不同语境下含义剧烈变化，说明该语义流形曲率极大。
\end{itemize}



\smalltitle{同调群 (Homology Group) — $H_n(X)$}

\begin{itemize}
\item   \textbf{数学定义}：拓扑空间中“洞”的数量与维度的代数度量（贝蒂数）。

\item   \textbf{MSC 对应}：\textbf{系统的高阶结构特征}。

\item   $H_0$：连通分量（有多少个独立的物体）。

\item   $H_1$：环/孔洞（是否存在逻辑循环或反馈回路）。

\item   $H_2$：空腔（是否存在被包裹的内部空间/自我）。
\end{itemize}
